\chapter{Conclusion}
\label{ch:conclusion}

A direct comparison of the co-digestion and \glsauto{cgl} only phases reveals that the co-digestion of \glsauto{fw} and \glsauto{cgl} resulted in higher biogas yields and more stable operation. The pH decreased faster with the co-digestion, while with \glsauto{cgl} only phase it did not decrease that fast, and even grew with the \glsauto{khco3} that was added to maintain the pH. These results support the hypothesis that \glsauto{cgl}, when co-digested with \glsauto{fw}, can enhance biogas production by balancing the \glsauto{cn_ratio} and providing readily available carbon, whereas \glsauto{cgl} alone may work, but would need a different setup and ajustments.

The \glsauto{cgl} only phase, should not be considered strictly worse than the co-digestion of \glsauto{fw} and \glsauto{cgl}, since both phases were performed in sequence, in the same reactor, and this could have given the \glsauto{cgl} phase disadvantages in a contaminated or exhausted environment.

Using \glsauto{cgl}, with some \glsauto{pgl} to make the necessary adjustment to balance the \glsauto{cn_ratio}, could be a viable solution for the \glsauto{cgl} problem. Its only major problem is that \glsauto{cgl} contains a lot of impurities that can affect the anaerobic digestion, and for a system that can reliably consume it, it will need a better control of its substrate.

During the whole process, the biggest production of \glsauto{ch4} was when \glsauto{vfa_alk} was \approx\num{1} with a pH of \approx\num{6.5}, during the \glsauto{cgl} only phase, even with not optimal values of pH and \glsauto{vfa_alk}, it was able to maintain a reasonable amount of production of \glsauto{ch4}, showing that it can be a suitable substrate for \glsauto{ad}.

The \glsauto{ch4} yield obtained (\qty{28.96}{\nml\CH\per\gram\VSfed}) is lower than values reported in the literature for similar co-digestion systems. The highest volumetric productivity (\qty{0.0446}{\litre\per\reactorlitre\per\day}) was observed during the co-digestion phase, while the highest specific yield (\qty{45.64}{\nml\CH\per\gram\VSfed}) occurred during the adaptation phase. This shows the importance of controlling the \glsauto{olr}, \glsauto{vfa_alk} and the amount of \glsauto{cgl} used to enhance the \glsauto{ch4} production.

Future studies could involve on how to better use the \glsauto{cgl} as a feed, to mitigate its downside, related to impurities and need for an extra carbon rich substrate to balance the feed.
